\begin{recipe}{Brood}
\label{recipe:brood}
\kicker[Brood,Basisrecept]{Brood \& basis · basisrecept}
\index[register]{Brood}
\index[register]{Broodmix}

\meta{Rijzen ca. 1 uur 15 \dvd 1 brood \dvd Oven 220\,°C}

\lettrine{V}{oor} dit brood gebruiken we een broodmix, Soezie van de Boerenbond
werkt bij ons het beste. Zelf voeg je alleen nog zout, gist, boter en water
toe.

\blockrule{Ingrediënten}
\begin{ingredients}
  \ing{500 g}{broodmix (bijv. Soezie van de Boerenbond)}\index[register]{Broodmix}
  \ing{280 ml}{water}
  \ing{15 g}{boter}
  \ing{8 g}{zout}
  \ing{8 g}{droge gist}
  \ingb{maanzaad, sesamzaad of zonnebloempitten, optioneel}
\end{ingredients}

\blockrule{Bereiding}
\tip{Onderschat het kneden niet: te kort kneden is echt niet goed. Is het
koud in huis? Zet het deeg dan te rijzen in de oven op de warmhoudstand van
30\,°C.}
\begin{steps}
  \item Doe de broodmix, de boter, het zout en de gist in een kom. Houd het
        zout en de gist daarbij uit elkaar, zout remt de werking van de gist.
  \item Voeg het water toe en kneed tot een soepel deeg. Doe de vliesjestest:
        kun je een stukje deeg uitrekken tot een dun vliesje zonder dat het
        scheurt? Dan is het genoeg. Reken op ongeveer 20 minuten, maar de
        test is betrouwbaarder dan de klok.
  \item Dek af met een natte theedoek en laat 30 à 40 minuten rijzen.
  \item Druk het deeg plat en geef het de gewenste vorm. Leg het in een
        ingevette bakvorm.
  \item Laat nog 40 minuten rijzen onder een natte theedoek.
  \item Bestrooi de bovenkant eventueel met maanzaad, sesamzaad of
        zonnebloempitten voor het bakken.
  \item Zet vlak voordat het brood de oven in gaat een bakje kokend water
        onderin de oven, of besproei de oven kort met water: stoom in de
        eerste 10 minuten geeft een dunnere, knapperigere korst en meer
        ovenrijs. Bak 30 à 35 minuten in een voorverwarmde oven op 220\,°C.
        Klop op de onderkant van het brood: klinkt het hol, dan is het gaar.
  \item Laat het brood minstens 20 minuten op een rooster afkoelen voor je
        het aansnijdt: vers uit de oven zet de kruim nog niet en snijd je
        het brood plat.
\end{steps}
\end{recipe}
